\documentclass[10pt,a4paper]{amsart}
\usepackage{amsmath,amssymb,amsthm,graphicx}
% Retained preamble for the Stacks Project chapter morphisms.tex (Morphisms of Schemes).
% The chapter file pinned-original-source.tex carries only \input{preamble} on its first line, so no
% preamble text lives inside the chapter itself.  The Stacks Project uses one single shared preamble for
% every chapter; the snapshot stages it verbatim in sources/stacks-chow/preamble.tex.  The block below is
% copied mechanically, line for line, from that file's lines 14-18 (the Xy-pic setup) and lines 154-217
% (the native theorem declarations together with the mathematical macros).  Nothing was added, removed or
% redefined.  The wrapper class and the amsmath/amssymb/amsthm/graphicx base packages are supplied by the
% assembler, so the class/\externaldocument part of the shared preamble is not repeated here.
% For commutative diagrams we use Xy-pic
\usepackage[all]{xy}

% We use 2cell for 2-commutative diagrams.
\xyoption{2cell}
% Theorem environments.
%
\theoremstyle{plain}
\newtheorem{theorem}[subsection]{Theorem}
\newtheorem{proposition}[subsection]{Proposition}
\newtheorem{lemma}[subsection]{Lemma}

\theoremstyle{definition}
\newtheorem{definition}[subsection]{Definition}
\newtheorem{example}[subsection]{Example}
\newtheorem{exercise}[subsection]{Exercise}
\newtheorem{situation}[subsection]{Situation}

\theoremstyle{remark}
\newtheorem{remark}[subsection]{Remark}
\newtheorem{remarks}[subsection]{Remarks}

\numberwithin{equation}{subsection}

% Macros
%
\def\lim{\mathop{\mathrm{lim}}\nolimits}
\def\colim{\mathop{\mathrm{colim}}\nolimits}
\def\Spec{\mathop{\mathrm{Spec}}}
\def\Hom{\mathop{\mathrm{Hom}}\nolimits}
\def\Ext{\mathop{\mathrm{Ext}}\nolimits}
\def\SheafHom{\mathop{\mathcal{H}\!\mathit{om}}\nolimits}
\def\SheafExt{\mathop{\mathcal{E}\!\mathit{xt}}\nolimits}
\def\Sch{\mathit{Sch}}
\def\Mor{\mathop{\mathrm{Mor}}\nolimits}
\def\Ob{\mathop{\mathrm{Ob}}\nolimits}
\def\Sh{\mathop{\mathit{Sh}}\nolimits}
\def\NL{\mathop{N\!L}\nolimits}
\def\CH{\mathop{\mathrm{CH}}\nolimits}
\def\proetale{{pro\text{-}\acute{e}tale}}
\def\etale{{\acute{e}tale}}
\def\QCoh{\mathit{QCoh}}
\def\Ker{\mathop{\mathrm{Ker}}}
\def\Im{\mathop{\mathrm{Im}}}
\def\Coker{\mathop{\mathrm{Coker}}}
\def\Coim{\mathop{\mathrm{Coim}}}

% Boxtimes
%
\DeclareMathSymbol{\boxtimes}{\mathbin}{AMSa}{"02}

%
% Macros for moduli stacks/spaces
%
\def\QCohstack{\mathcal{QC}\!\mathit{oh}}
\def\Cohstack{\mathcal{C}\!\mathit{oh}}
\def\Spacesstack{\mathcal{S}\!\mathit{paces}}
\def\Quotfunctor{\mathrm{Quot}}
\def\Hilbfunctor{\mathrm{Hilb}}
\def\Curvesstack{\mathcal{C}\!\mathit{urves}}
\def\Polarizedstack{\mathcal{P}\!\mathit{olarized}}
\def\Complexesstack{\mathcal{C}\!\mathit{omplexes}}
% \Pic is the operator that assigns to X its picard group, usage \Pic(X)
% \Picardstack_{X/B} denotes the Picard stack of X over B
% \Picardfunctor_{X/B} denotes the Picard functor of X over B
\def\Pic{\mathop{\mathrm{Pic}}\nolimits}
\def\Picardstack{\mathcal{P}\!\mathit{ic}}
\def\Picardfunctor{\mathrm{Pic}}
\def\Deformationcategory{\mathcal{D}\!\mathit{ef}}

% --- appended 2026-10-12 (Chief): the project own definitions of the reference/slogan/history
% display environments, copied mechanically from the shared Stacks preamble lines 8-12
% (\usepackage{verbatim} plus the three \newenvironment lines, which the chapter preamble keeps
% verbatim from the project preamble). Nothing added, removed or redefined: these environments
% carry no mathematical content and the project itself typesets them as comments. Without these
% three lines a counted statement whose own text uses them fails with Environment slogan undefined.
\usepackage{verbatim}
\newenvironment{reference}{\comment}{\endcomment}
\newenvironment{slogan}{\comment}{\endcomment}
\newenvironment{history}{\comment}{\endcomment}
\usepackage{cleveref}
\crefname{originaldefinition}{original definition}{original definitions}
\Crefname{originaldefinition}{Original definition}{Original definitions}
\makeatletter
\newlabel{definition-property-local}{{01SR}{source}{}{sourcepointer-1}{}}
\newlabel{definition-property-local@cref}{{[originaldefinition][0][]01SR}{[1][0][]source}{}{}{}}
\makeatother
\crefname{originallemma}{original lemma}{original lemmas}
\Crefname{originallemma}{Original lemma}{Original lemmas}
\makeatletter
\newlabel{lemma-locally-P}{{01ST}{source}{}{sourcepointer-2}{}}
\newlabel{lemma-locally-P@cref}{{[originallemma][0][]01ST}{[1][0][]source}{}{}{}}
\makeatother
\crefname{originallemma}{original lemma}{original lemmas}
\Crefname{originallemma}{Original lemma}{Original lemmas}
\makeatletter
\newlabel{lemma-locally-finite-type-characterize}{{01T2}{source}{}{sourcepointer-3}{}}
\newlabel{lemma-locally-finite-type-characterize@cref}{{[originallemma][0][]01T2}{[1][0][]source}{}{}{}}
\makeatother
\crefname{originallemma}{original lemma}{original lemmas}
\Crefname{originallemma}{Original lemma}{Original lemmas}
\makeatletter
\newlabel{lemma-quasi-finite-at-point-characterize}{{01TH}{source}{}{sourcepointer-4}{}}
\newlabel{lemma-quasi-finite-at-point-characterize@cref}{{[originallemma][0][]01TH}{[1][0][]source}{}{}{}}
\makeatother
\title{Local quasi-finiteness is a local property of ring maps and is characterized by it on affine charts}
\author{The Stacks Project Authors}
\date{stacks-project:morphisms.tex:local-snapshot-20261012; GNU FDL 1.2 or later}
\begin{document}
\maketitle
\noindent\textit{Editorial note.} All delivered mathematical text is the author's own, copied from the pinned Stacks Project chapter file pinned-original-source.tex; no statement or proof step was rewritten, reordered or removed.\par\medskip
\noindent\textit{Editorial note.} Following the multi-claim convention for this source, source\_url is the source-level URL https://stacks.math.columbia.edu/ and the source label lemma-locally-quasi-finite-characterize is carried in claim\_key and here rather than in a per-claim tag URL.\par\medskip
\noindent\textit{Editorial note.} The preamble is the chapter's shared Stacks Project preamble, retained as sources/stacks-morphisms/retained-preamble.tex: a line-for-line excerpt of the shared Stacks preamble staged with this snapshot. No macro was invented or redefined, and this card adds no preamble text of its own.\par\medskip
\noindent\textit{Editorial note.} Statement span 3838-3859 and proof span 3861-3894 come from the halfA pool card operations/pool-cards/stacks-morphisms-halfA.json and were re-read in the pinned chapter file before this card was written.\par\medskip
\noindent\textit{Editorial note.} The statement and proof contain several counts of conditions by letters (a), (b), (c) referring to the cited local-property definition; these letters are the author's and are kept verbatim.\par\medskip
\noindent\textit{Editorial note.} 4 cross-references occur inside the counted spans; they are listed in external\_references as navigation pointers only.\par\medskip
\noindent\textit{Allowed standard prerequisites (editorial).}\begin{itemize}
\item The definition of a local property of ring maps, with its three conditions (a), (b), (c) on localizations and on covers of the ring. Cited as a same-chapter definition at proof lines 4079-4081 and used throughout the first paragraph. It is the abstract notion whose conditions the author verifies for quasi-finiteness, and which then lets the same-chapter lemma on local properties be applied.
\item The same-chapter statement that a local property of ring maps gives the corresponding local property of morphisms, applied with P = quasi-finite. Cited at proof lines 4102-4104. Reduces the four equivalences of the statement to the single algebraic property.
\item The same-chapter characterization of quasi-finiteness at a point by the corresponding property of the local ring map. Cited at proof lines 4107-4110. Identifies the morphism-level property with the ring-level property P, completing the reduction.
\item That being of finite type is already known to be a local property of ring maps, together with the fact that quasi-finiteness at a prime depends only on the local ring. Cited at proof lines 4082-4090. Supplies conditions (a) and (b) for quasi-finiteness and the local-ring reduction used in condition (c).
\end{itemize}\medskip
\ifdefined\hypertarget\hypertarget{sourcepointer-1}{}\fi
\noindent\textit{Source reference.} Original definition 01SR: focused-seeds/AGP/sources/stacks-snapshot-20261011/tags.txt, source lines 2348--2348. Same-chapter reference used in the retained author text at source line(s) 2471, 2480, 2483, 2490, 2699, 3867, 3875, 4122, 5812, 7074, 7695, 8116. It is supplied only as a navigation pointer so that the retained text compiles; the counted item delivers the author's own statement and the author's complete proof, and no part of the referenced item is copied in or substituted for the counted argument.\par\medskip
\ifdefined\hypertarget\hypertarget{sourcepointer-2}{}\fi
\noindent\textit{Source reference.} Original lemma 01ST: focused-seeds/AGP/sources/stacks-snapshot-20261011/tags.txt, source lines 2350--2350. Same-chapter reference used in the retained author text at source line(s) 2518, 2696, 3887, 5809, 7071, 7692, 8113. It is supplied only as a navigation pointer so that the retained text compiles; the counted item delivers the author's own statement and the author's complete proof, and no part of the referenced item is copied in or substituted for the counted argument.\par\medskip
\ifdefined\hypertarget\hypertarget{sourcepointer-3}{}\fi
\noindent\textit{Source reference.} Original lemma 01T2: focused-seeds/AGP/sources/stacks-snapshot-20261011/tags.txt, source lines 2359--2359. Same-chapter reference used in the retained author text at source line(s) 2717, 2739, 2801, 2863, 3019, 3035, 3157, 3348, 3613, 3707, 3868, 5197, 5282, 6702, 9325, 9390, 9440, 11421, 13422, 13686. It is supplied only as a navigation pointer so that the retained text compiles; the counted item delivers the author's own statement and the author's complete proof, and no part of the referenced item is copied in or substituted for the counted argument.\par\medskip
\ifdefined\hypertarget\hypertarget{sourcepointer-4}{}\fi
\noindent\textit{Source reference.} Original lemma 01TH: focused-seeds/AGP/sources/stacks-snapshot-20261011/tags.txt, source lines 2374--2374. Same-chapter reference used in the retained author text at source line(s) 3749, 3770, 3781, 3797, 3803, 3893, 3973, 3988, 4037, 5657, 5667, 5674, 7816, 11157, 13319, 13323, 15193, 15510. It is supplied only as a navigation pointer so that the retained text compiles; the counted item delivers the author's own statement and the author's complete proof, and no part of the referenced item is copied in or substituted for the counted argument.\par\medskip
\begin{lemma}
\label{lemma-locally-quasi-finite-characterize}
Let $f : X \to S$ be a morphism of schemes.
The following are equivalent
\begin{enumerate}
\item The morphism $f$ is locally quasi-finite.
\item For every pair of affine opens $U \subset X$, $V \subset S$
with $f(U) \subset V$ the ring map
$\mathcal{O}_S(V) \to \mathcal{O}_X(U)$ is quasi-finite.
\item There exists an open covering $S = \bigcup_{j \in J} V_j$
and open coverings $f^{-1}(V_j) = \bigcup_{i \in I_j} U_i$ such
that each of the morphisms $U_i \to V_j$, $j\in J, i\in I_j$
is locally quasi-finite.
\item There exists an affine open covering $S = \bigcup_{j \in J} V_j$
and affine open coverings $f^{-1}(V_j) = \bigcup_{i \in I_j} U_i$ such
that the ring map $\mathcal{O}_S(V_j) \to \mathcal{O}_X(U_i)$ is
quasi-finite, for all $j\in J, i\in I_j$.
\end{enumerate}
Moreover, if $f$ is locally quasi-finite then for
any open subschemes $U \subset X$, $V \subset S$ with $f(U) \subset V$
the restriction $f|_U : U \to V$ is locally quasi-finite.
\end{lemma}

\begin{proof}
For a ring map $R \to A$ let us define
$P(R \to A)$ to mean ``$R \to A$ is quasi-finite''
(see remark above lemma).
We claim that $P$ is a local property of ring maps.
We check conditions (a), (b) and (c) of Definition
\ref{definition-property-local}. In the proof of
Lemma \ref{lemma-locally-finite-type-characterize}
we have seen that (a), (b) and (c) hold for the property
of being ``of finite type''. Note that, for a finite type ring map
$R \to A$, the property $R \to A$ is quasi-finite at $\mathfrak q$
depends only on the local ring $A_{\mathfrak q}$ as an
algebra over $R_{\mathfrak p}$ where $\mathfrak p = R \cap \mathfrak q$
(usual abuse of notation). Using these remarks (a), (b) and (c) of
Definition \ref{definition-property-local} follow immediately.
For example, suppose $R \to A$ is a ring map
such that all of the ring maps $R \to A_{a_i}$ are quasi-finite
for $a_1, \ldots, a_n \in A$ generating the unit ideal.
We conclude that $R \to A$ is of finite type. Also, for any
prime $\mathfrak q \subset A$ the local ring $A_{\mathfrak q}$
is isomorphic as an $R$-algebra to the local ring
$(A_{a_i})_{\mathfrak q_i}$ for some $i$ and some
$\mathfrak q_i \subset A_{a_i}$. Hence we conclude that
$R \to A$ is quasi-finite at $\mathfrak q$.

\medskip\noindent
We conclude that Lemma \ref{lemma-locally-P} applies with $P$
as in the previous paragraph.
Hence it suffices to prove that $f$ is locally quasi-finite is
equivalent to $f$ is locally of type $P$. Since $P(R \to A)$
is ``$R \to A$ is quasi-finite'' which means $R \to A$ is
quasi-finite at every prime of $A$, this follows from
Lemma \ref{lemma-quasi-finite-at-point-characterize}.
\end{proof}

\end{document}
